\section{The actual adelic projectors and one common arithmetic row family}
\label{sec:arithmetic-rows-v46}

This chapter constructs the row object on which the retained-weight
and cylinder-count arguments operate.  The input is one original
datum \(D\) and one original adelic \(\mathrm{SG.HidaSetting}(D)\),
denoted by \(X\).  This is the setting of lattices, adelic measures,
discrete-series representations and row operators.  It is a different
record from the ordinary Hecke-algebra setting used in the preceding
Hida chapters.  Nothing here identifies their objects merely because
both records have the word Hida in their name.

Write \(n=D.R.r\), \(f=D.R.f\), \(s=D.R.rS=D.R.m/2-1\),
\(k_0=D.R.k\), and \(q_j=D.\mathrm{host}.g(j)\), for \(j=0,1\).
Thus \(n=2s+k_0+1\), and \(q_0,q_1\) are distinct primes splitting
completely in the original CM field.  Transport through \(X.\iota_E\)
is understood throughout.  Put \(g=q_0q_1\), let \(K\) be the
original finite-adelic lattice stabilizer, and let
\[
 N=K(g)=\{x\in K:(x-1)\widehat L\subseteq g\widehat L\},
 \qquad
 \mathcal G=\prod_{j=0}^1\prod_{v\in D.\mathrm{nd}.\mathrm{PlG}(j)}
                  \mathrm{GL}_n(\mathbb F_{q_j}).
 \tag{ARO.1}
\]
Every measure, fundamental domain, infinity representation, frame,
and lattice in this chapter is the corresponding object of \(X\).
In particular \(\mu=X.B.\mu_f\) is never renormalized.  Let
\(V=\mu(N)\), with finite measures identified with their real values.
The finite labels are the actual monic degree-\(n\) polynomials with
nonzero constant coefficient; \(E_{j,P}=D.\mathrm{ntypes}.eSer(j,P)\)
are their original central orthogonal idempotents.  An invalid
polynomial has idempotent zero.  The normalized-type record, the
given setting \(X\), and, where expressly used below, the original
\(\mathrm{RowsFinDim}\), \(\mathrm{GapHyp}\), \(\mathrm{AutTie}\),
\(\mathrm{RowDictTie}\), and \(\mathrm{LabTie}\) are retained inputs.
Their simultaneous existence in a chosen global family is not assumed
to follow from this chapter.

\subsection{Reduction, the level character, and the exact Haar factor}

% BEGIN ARO1
\paragraph{ARO1. Reduction on the original lattice and its exact kernel.}
\label{claim:ARO1}
Choose once, for each selected place \(w=w(j,v)\), a basis of
\(L/wL\), and lift it to \(L\).  The action on these quotients defines
a homomorphism \(\rho:K\to\mathcal G\) with
\[
                         \ker\rho=N.                 \tag{ARO.2}
\]
The same assertion for one auxiliary prime says that the kernel of
the corresponding factor of \(\rho\) is \(K(q_j)\).
\begin{proof}
The localization of a full lattice at a finite prime is a free module
of rank \(n\) over its discrete valuation ring.  Reduction modulo its
maximal ideal has dimension \(n\), so the asserted bases exist.
The ordinary lattice maps onto this reduction: a localized
denominator is invertible modulo the maximal ideal and can be
replaced by its inverse in the residue field.  Since that field is
\(\mathbb F_{q_j}\), its coefficients can also be lifted by integers.
Every \(x\in K\) and its inverse preserve the completed lattice;
their reduced matrices are inverse.  Reducing products proves the
homomorphism assertion.  A different lifted basis conjugates the
reduced matrix and does not affect whether it is the identity.

Here are the two points needed to identify the kernel on the whole
adelic lattice.  At a selected split prime, the original frame and
dual frame have pairing matrix the identity modulo that prime.
This follows directly from \(\mathrm{gramIntegral}\),
\(\mathrm{formsAtZero}\), \(\mathrm{bw\_dual}\), and the
complementary residues of \(\mathrm{ew}\) and its conjugate.
Indeed an idempotent in a residue field is zero or one;
\(\mathrm{ew\_orth}\) and \(\mathrm{ew\_full}\) make the two
residues complementary.  Pairing a linear relation with the dual
vectors proves independence, and \(\mathrm{latticeSpan}\) proves
spanning.  Thus the pairing of the two conjugate local lattices is
perfect.  Unitarity identifies their action matrices as inverse
transposes in dual bases.  Identity reduction at \(w\) therefore
implies identity reduction at its conjugate \(\bar w\).

The selected places and their conjugates exhaust all primes above
\(q_0q_1\).  Complete splitting means these primes are distinct and
occur to exponent one.  In each local free lattice, zero reduction
of \((x-1)y\) means membership in its maximal-ideal multiple.
At the remaining primes \(g\) is a unit.  Consequently identity
reduction in every selected factor is equivalent to
\((x-1)\widehat L\subseteq g\widehat L\), component by component.
This proves (ARO.2); the proof with only the primes above \(q_j\)
gives the last assertion.  The argument uses the original completed
lattice, so it does not replace adelic congruence by congruence of a
chosen collection of rational vectors.
\end{proof}
% END ARO1

% BEGIN ARO2
\paragraph{ARO2. Surjectivity of the actual reduction map.}
\label{claim:ARO2}
The homomorphism of ARO1 is surjective.  Hence
\(K/N\simeq\mathcal G\), with exactly \(f\) factors above each
auxiliary prime.
\begin{proof}
At a completely split finite place the CM algebra is a product of
two copies of the base local field, interchanged by conjugation.
The unitary equation determines the action on the second factor
from its action on the first.  In the dual local bases supplied in
ARO1 it reads \(A'=(A^{-1})^{\mathsf t}\).  The lattice stabilizer
therefore identifies with \(\mathrm{GL}_n(\mathcal O_v)\).
An arbitrary matrix in \(\mathrm{GL}_n(\mathbb F_{q_j})\) lifts by
lifting its entries: its determinant remains a unit, so its inverse
also has integral entries.  This constructs an actual local unitary
lattice automorphism with the prescribed reduction at that place.

Do this at the finitely many selected places and take the identity
at every other finite place.  The resulting tuple is a finite adele,
satisfies the unitary equation at every place and preserves every
local lattice.  It is therefore an element of the original \(K\).
Changing from these convenient local bases to the bases defining
\(\rho\) only conjugates each target matrix; prescribe the conjugated
matrix when making the lift.  All factors can be prescribed
independently.  This proves surjectivity without a rational strong
approximation assertion.  ARO1 now gives the quotient isomorphism.
The number of factors is the original equality
\(|D.\mathrm{nd}.\mathrm{PlG}(j)|=f\).
\end{proof}
% END ARO2

% BEGIN ARO3
\paragraph{ARO3. Haar volume on each reduction fiber.}
\label{claim:ARO3}
One has \(0<V<\infty\).  Each fiber of \(\rho\) has measure \(V\).
If \(H\leq\mathcal G\), then
\[
 \mu(\rho^{-1}H)=|H|V,
 \qquad \mu(K)=|\mathcal G|V.                         \tag{ARO.3}
\]
\begin{proof}
The congruence subgroup \(N\) is open and contains the identity.
It is a closed subgroup of the compact lattice stabilizer \(K\).
A nonzero Haar measure is positive on every nonempty open set and
finite on compact sets, proving the inequalities.  By ARO1--ARO2
a reduction fiber is one coset \(xN\); left invariance gives measure
\(V\).  The fibers are disjoint measurable sets.  Finite additivity
over the fibers in \(H\) gives (ARO.3).  In particular no counting
measure of total mass one has been substituted for \(\mu\).
\end{proof}
% END ARO3

% BEGIN ARO4
\paragraph{ARO4. Inflation is independent of the reduction basis.}
\label{claim:ARO4}
For a central element \(a\in\mathbb C[\mathcal G]\), the original
function \(\mathrm{egA}(X,a)\) is exactly
\[
 \operatorname{inf}(a)(x)=
 \begin{cases}a(\rho(x))/V,&x\in K,\\0,&x\notin K.\end{cases}
 \tag{ARO.4}
\]
It is locally constant, compactly supported and bi-invariant under
\(N\).  For the original finite event idempotent \(E_S\), this
applies simultaneously to every set of tuples \(S\).
\begin{proof}
The original definition selects any tuple of reduction matrices in
any admissible lifted bases.  Such bases really are bases: their
spanning condition on a rank-\(n\) quotient of dimension \(n\)
implies independence.  Two choices are related by an invertible
change of basis in each factor, so their tuples are conjugate in
\(\mathcal G\).  Coefficients of a central group-algebra element are
constant under conjugacy.  ARO1 supplies one tuple for each \(x\in K\),
and the original denominator is precisely \(V\).  This proves
(ARO.4), including the zero value outside \(K\).
Finite fibers are open by ARO1; their union is \(K\), and \(K\)
is also closed.  This proves local constancy and compact support.
Multiplication on either side by \(N=\ker\rho\) changes no value.
RRT12 proves centrality of every actual \(E_S\), so the conclusion
holds for all these events with the same \(\rho\) and \(V\).
\end{proof}
% END ARO4

% BEGIN ARO5
\paragraph{ARO5. The exact finite convolution and involution comparison.}
\label{claim:ARO5}
For a function \(h\) supported in \(K\) and bi-invariant under
\(N\), define \(\widehat h(\gamma)=Vh(x_\gamma)\), where
\(\rho(x_\gamma)=\gamma\).  Then this definition is independent of
the representatives, is injective on such functions, and satisfies
\[
 \widehat{h_1*h_2}=\widehat h_1\widehat h_2,
 \qquad
 \widehat{h^*}=\widehat h^{\,*},
 \qquad h^*(x)=\overline{h(x^{-1})}.
 \tag{ARO.5}
\]
\begin{proof}
Representatives of a fiber differ by \(N\), by ARO1.
The values of \(h\) on \(K\) can be recovered by dividing
\(\widehat h\) by the positive number \(V\), proving injectivity.
For convolution partition its domain of integration into reduction
fibers.  At \(x_\gamma\) the integrand on \(x_\delta N\) is the
constant \(h_1(x_\delta)h_2(x_\delta^{-1}x_\gamma)\).  ARO3 gives
\[
 (h_1*h_2)(x_\gamma)
 =V\sum_{\delta\in\mathcal G}
           h_1(x_\delta)h_2(x_{\delta^{-1}\gamma}).
\]
Multiplication by the outer \(V\) gives the group-algebra product;
there are two factors \(V\), one from the integral and one from the
transform.  If \(x\notin K\), both nonzero factors in the integrand
would force \(x\in K\), so the convolution is zero.  The convolution
is again bi-invariant.  Inversion sends a fiber to the inverse
fiber, and \(V\) is real, proving the involution formula.  All
integrals here are finite sums of bounded functions on compact sets.
By ARO4, \(\widehat{\operatorname{inf}(a)}=a\).
\end{proof}
% END ARO5

% BEGIN ARO6
\paragraph{ARO6. The finite event algebra with its invalid-label convention.}
\label{claim:ARO6}
Let \(\mathcal T\) be the finite set of valid two-prime tuples.
For any set \(S\) of polynomial tuples, the literal \(E_S\) satisfies
\[
 E_SE_T=E_{S\cap T},\quad E_S^*=E_S,\quad E_{\mathrm{all}}=1,
 \quad E_S=\sum_{t\in\mathcal T\cap S}E_{\{t\}}.
 \tag{ARO.6}
\]
In particular arbitrary sets cause no infinite-sum issue.
\begin{proof}
A monic polynomial of degree \(n\) is uniquely determined by its
\(n\) coefficients below the leading coefficient.  Conversely each
coefficient tuple gives that polynomial.  The original definition of
\(E_S\) is consequently the sum of tensor products of the local
\(E_{j,P}\), indexed once by each monic tuple in \(S\).  Invalid
labels have a zero factor.  Removing them leaves exactly the finite
sum in (ARO.6).  RRT12 gives the product and total-sum identities.

For clarity self-adjointness follows even if it is not included as
a separate field of the normalized-type record.  In the left regular
representation a central element acts by a scalar on each irreducible
summand.  A central idempotent has scalar zero or one.  The regular
representation is unitary for its standard finite counting inner
product and faithful, so this idempotent equals its adjoint.
Tensor products and finite sums preserve that equality.  This proves
the involution assertion for all original event projectors.
\end{proof}
% END ARO6

% BEGIN ARO7
\paragraph{ARO7. The original frame gives dual bases at both orientations.}
\label{claim:ARO7}
At a selected place let \(b\in\{0,1\}\) be the residue of
\(X.\mathrm{ew}\).  A residue basis, in four consecutive blocks,
is
\[
 (e_0,\ldots,e_{s-1};\ c_0,\ldots,c_{k_0-1};\ z;
                         f_{s-1},\ldots,f_0),        \tag{ARO.7}
\]
where \(c_i=\beta_i\) if \(b=1\), and
\(c_i=\beta^d_{k_0-1-i}\) if \(b=0\).
Its dual frame is obtained by replacing \(e\) by \(f\),
\(z\) by \(-z\), and \(c\) by the complementary \(\beta\)-frame
with the same reversal.  These statements concern the original
frame vectors modulo the place.
\begin{proof}
The complementary-residue argument in ARO1 gives
\(\mathrm{ew}\equiv b\) and \(\overline{\mathrm{ew}}\equiv1-b\).
The identities \(\beta_i=\mathrm{ew}\,\beta_i\) and
\(\beta_i^d=\overline{\mathrm{ew}}\,\beta_i^d\) make one of the
two \(\beta\)-frames zero at this place and the other the active
frame.  The congruences in the setting give
\(h(e_a,f_c)=\delta_{ac}\), \(h(z,z)=-1\), and the corresponding
active dual \(\beta\)-pairings \(\delta_{ij}\); all off-block
pairings vanish.  Pairing (ARO.7) with its displayed dual therefore
gives the identity matrix.  Independence follows by pairing any
linear relation with the dual frame.  The original lattice-spanning
condition makes the active frame span, so it is a basis, not just an
independent subspace.  Reversal in the zero orientation is necessary
because conjugation interchanges the active and dual flags.
The changes from the basis of ARO1 to these bases are invertible
matrices fixed before the group element is varied.
\end{proof}
% END ARO7

% BEGIN ARO8
\paragraph{ARO8. Exact reduction of the concrete level subgroup.}
\label{claim:ARO8}
Let \(J_q\) be the subgroup of upper unitriangular
\(n\times n\) matrices over \(\mathbb F_q\) satisfying
\[
                  M_{a,s+k_0}=0\quad(s\leq a<s+k_0).
 \tag{ARO.8}
\]
There are matrices \(P_{j,v}\), fixed before \(x\), such that,
for every \(x\in K\),
\[
 x\in X.K_2\quad\Longleftrightarrow\quad
 P_{j,v}^{-1}\rho(x)_{j,v}P_{j,v}\in J_{q_j}
                    \quad\hbox{for every }j,v.       \tag{ARO.9}
\]
\begin{proof}
In the basis of ARO7 the defining level congruences say that
\((x-1)e_a\) is in the span of earlier \(e\)'s,
\((x-1)z\) is in their full span, and
\((x-1)\beta_i\) is in that span together with the earlier
\(\beta\)'s.  Thus, in the orientation \(b=1\), the \(e\),
active \(\beta\), and \(z\) columns have zero entries below the
diagonal, diagonal entry one, and the extra zero entries (ARO.8).
Applying the conjugate-place conditions and the identity
\(h(u,xv)=h(x^{-1}u,v)\) supplies the same conditions on the
reversed \(f\)-block.  In the orientation \(b=0\), the active
\(\beta^d\)-block is controlled by those conjugate inverse
conditions; the reversal in ARO7 turns them into upper-triangular
conditions.  The \(e\) and \(z\) columns are controlled directly.
For the \(z\) column the coefficients in the active middle block
are zero in both orientations.  These are exactly all entries below
the diagonal, all diagonal entries, and (ARO.8).

The converse can be checked without a volume-index argument.
Upper unitriangular matrices satisfying (ARO.8) form a subgroup:
the only possible intermediate terms in the forbidden column are
other forbidden entries.  Hence their inverses have the same shape.
The displayed shape gives the required flag inclusions for the
\(e\) and \(z\) vectors directly.  For \(b=1\) it gives them for
\(\beta\) directly; for \(b=0\) those vectors vanish at this
place and the conjugate-place inverse matrix supplies their
conditions.  Applying this at both places in every conjugate pair
proves the original inclusions modulo every prime above \(g\).
Chinese remaindering for the completed lattice then gives the
original inclusions modulo \(g\), with their original spans.
Outside these primes \(g\) is a unit and imposes no condition.

Finally multiply the basis by a diagonal matrix equal to one on
the \(e\) and \(z\) blocks, to \((-1)^i\) on the reversed
\(f\)-block, and also to \((-1)^i\) on the middle block when
\(b=0\).  This preserves the zero pattern and diagonal entries.
It is the sign adjustment used for the character below.  Compose
it with the fixed basis change of ARO7 to obtain \(P_{j,v}\).
This proves both directions of (ARO.9) on the original subgroup.
\end{proof}
% END ARO8

% BEGIN ARO9
\paragraph{ARO9. Factorization and nontriviality of the actual level character.}
\label{claim:ARO9}
For each selected place there is a nontrivial additive character
\(\psi_{j,v}:\mathbb F_{q_j}\to\mathbb C^\times\) such that, on
the subgroup of ARO8,
\[
 X.\chi_2(x)=\prod_{j,v}\psi_{j,v}
   \bigl(\phi_J(P_{j,v}^{-1}\rho(x)_{j,v}P_{j,v})\bigr).
 \tag{ARO.10}
\]
Here \(\phi_J\) is exactly the original sum of entries: the
successive entries in the \(e\)-chain, the two links through \(z\),
the successive entries in the reversed \(f\)-chain, and the
successive entries in the middle block.  Empty chains contribute zero.
\begin{proof}
Use Chinese remaindering in \(\mathcal O_E/g\mathcal O_E\) to
choose \(\varepsilon_{j,v}\) equal to one at \(w(j,v)\) and zero
at every other prime above \(g\), including its conjugate.
It is divisible by \(q_{1-j}\).  Define
\[
 \psi_{j,v}(a)=\psi_f(\varepsilon_{j,v}\widetilde a/g),
 \tag{ARO.11}
\]
where \(\widetilde a\) is any integer lift.  Changing the lift
adds an integral adele because \(q_{1-j}\mid\varepsilon_{j,v}\).
The given additive character is trivial on such adeles, proving
well-definedness and additivity.  Its value on the conjugate CRT
idempotent is the same, by the original rational/infinity character
identity and invariance of the real trace under CM conjugation.

Nontriviality can be checked by a trace calculation in the original
field.  Put \(a=\varepsilon_{j,v}/q_{1-j}\in\mathcal O_E\).
Modulo \(q_j\), its absolute trace is the sum of its residues at
all primes above \(q_j\), since the integral algebra there is the
product of copies of \(\mathbb Z_{q_j}\).  This sum is
\(q_{1-j}^{-1}\neq0\) in \(\mathbb F_{q_j}\).  The rational
character identity therefore gives
\(\psi_{j,v}(1)=\exp(-2\pi i\operatorname{Tr}_{E/\mathbb Q}(a)/q_j)
\neq1\).  This uses complete splitting, already part of the
original datum, and does not posit a new conductor normalization.

Let \(C(x)\) be the original argument of \(\chi_2\), namely
\[
 C(x)=\sum_a h(f_a,x\,\operatorname{next}(e,z)_a)
                 +\sum_{i<k_0-1}h(\beta_i^d,x\beta_{i+1}).
\]
It is integral because all frame vectors and their transforms are
in the integral lattice.  In the frame of ARO7 its residue at
\(w\) is the sum of the direct \(e\)-chain entries and, in
orientation one, the middle-chain entries.  At \(\bar w\), use
\(\overline{h(u,xv)}=h(v,x^{-1}u)\).  This gives the reversed
\(f\)-chain entries of the inverse matrix, a minus sign at the
\(z\)-link because the dual of \(z\) is \(-z\), and, in
orientation zero, the reversed middle-chain entries of the inverse.
For upper unitriangular matrices the first superdiagonal of the
inverse is the negative of that of the matrix.  The diagonal signs
at the end of ARO8 change these negatives into the positive
successive entries in \(\phi_J\).  The two links crossing the
middle block have no extra product term, precisely because the
forbidden column in (ARO.8) is zero.  Thus the sum of the two
conjugate residues of \(C(x)\) is the stated \(\phi_J\).

Subtract from \(C(x)\) the sum of its residues times the CRT
idempotents and their conjugates.  The difference lies in
\(g\widehat{\mathcal O}_E\).  Divide by \(g\) and apply \(\psi_f\).
The difference contributes one, and additivity, (ARO.11), and the
equal conjugate characters give (ARO.10).  Every sign and every
local factor is therefore tied to the original character.
\end{proof}
% END ARO9

% BEGIN ARO10
\paragraph{ARO10. The level volume and its finite character idempotent.}
\label{claim:ARO10}
Put \(c_j=|J_{q_j}|\), and let
\[
 a_{j,v}=P_{j,v}e_J(q_j,n,k_0,s,\psi_{j,v})P_{j,v}^{-1}
                 \quad\hbox{in }\mathbb C[\mathrm{GL}_n(\mathbb F_{q_j})].
\]
Then, for the original function \(e=X.B.e\),
\[
 \mu(X.K_2)=c_0^f c_1^f V,
 \qquad \widehat e=\bigotimes_{j,v}a_{j,v},
 \qquad e(1)=\frac1{c_0^f c_1^f V}>0.                 \tag{ARO.12}
\]
\begin{proof}
By ARO2 and ARO8 the reduction image of \(X.K_2\) is the full
product of the conjugated subgroups \(J_{q_j}\), and its kernel is
\(N\).  ARO3 gives the volume formula.  The identity belongs to
each \(J_q\), so \(c_j>0\).  The formula defining the original
\(e\) is \(\overline{\chi_2}/\mu(X.K_2)\) on \(X.K_2\),
and zero off it.  ARO9 expresses this character as the product of
the local \(\psi_{j,v}\circ\phi_J\).  Its conjugate is its
inverse, because finite additive characters have absolute value
one.  Multiplying by \(V\) and using the volume formula gives
exactly the normalized subgroup-character idempotents in (ARO.12).
Their values at the identity are \(1/c_j\), proving the last
formula.  The tensor equality is a coefficient equality at every
element, including those outside the level image.
\end{proof}
% END ARO10

% BEGIN ARO11
\paragraph{ARO11. Independence of the chosen nontrivial finite character.}
\label{claim:ARO11}
If \(Z\) is central in \(\mathbb C[\mathrm{GL}_n(\mathbb F_q)]\)
and \(\psi,\psi'\) are nontrivial additive characters, then
\[
 [1](e_J(\psi)Z)=[1](e_J(\psi')Z),
 \qquad [1]e_J(\psi)=[1]e_J(\psi')=1/|J_q|.
 \tag{ARO.13}
\]
\begin{proof}
The additive group of the prime field is cyclic of prime order.
Every nontrivial character is faithful and every character is
determined by its value at one.  Thus
\(\psi'(x)=\psi(cx)\) for some \(c\in\mathbb F_q^\times\).
Assign the following exponents to the four blocks of ARO7:
\[
 0,\ldots,s-1;\quad 0,\ldots,k_0-1;\quad s;\quad
                                  s+1,\ldots,2s.
\]
Let \(T\) be diagonal with entries \(c\) to these powers.
Every entry occurring in \(\phi_J\) has exponent difference one.
Therefore \(\phi_J(T^{-1}MT)=c\phi_J(M)\).  Diagonal conjugation
preserves both the upper unitriangular shape and its forbidden
zero column, so it permutes \(J_q\).  Reindexing the character
average shows that the two idempotents are conjugate by \(T\)
(with the direction reversed if the inverse convention is used).
The identity coefficient of a conjugate is unchanged, and \(Z\)
commutes with \(T\), proving the first equality.  At the identity
the character value is one, proving the second.  The displayed
exponents also cover \(s=0\) and an empty middle chain.
\end{proof}
% END ARO11

% BEGIN ARO12
\paragraph{ARO12. The original cylinder identity including its Haar normalization.}
\label{claim:ARO12}
Choose any nontrivial \(\psi_j^0\) on each \(\mathbb F_{q_j}\),
and put \(e_T(S)=e*\operatorname{inf}(E_S)\).  For every finite
set \(S\) of one-prime tuples,
\[
 \operatorname{Re}e_T(\operatorname{cyl}(j,S))(1)
   =e(1)\sum_{t\in S}\prod_v
      \frac{\operatorname{Re}[1](e_J(\psi_j^0)E_{j,t_v})}
           {\operatorname{Re}[1]e_J(\psi_j^0)}.
 \tag{ARO.14}
\]
This is the full original \(\mathrm{EtOne}\), on the given measure.
\begin{proof}
ARO4--ARO5 give \(\widehat{e_T(S)}=\widehat e E_S\).
ARO10 identifies \(\widehat e\) as the actual tensor product of
conjugated local character idempotents.  In the identity coefficient
of its product with a central label projector, the conjugating
matrix disappears.  ARO11 then replaces every local \(\psi_{j,v}\)
by the chosen \(\psi_j^0\).  Apply the finite tensor-cylinder
identity RRT12: all complementary-prime coordinates sum to their
full identity, leaving exactly the sum of products in (ARO.14)
times \([1]\widehat e\).  Invalid labels contribute zero by ARO6.

Each numerator is real and nonnegative.  Indeed the product of a
central self-adjoint idempotent and a subgroup-character idempotent
is a self-adjoint idempotent; its identity coefficient is
\([1](b^*b)=\sum_\gamma|b(\gamma)|^2\).  This is also the
positivity calculation of RRT2.  The denominators are \(1/c_j>0\).
Finally divide both finite identity coefficients by the same
\(V>0\), using ARO5 and ARO10.  This yields the literal adelic
formula with \(e(1)\), proving (ARO.14).  The choice of finite
characters changes neither the measure nor any projector \(e_T\).
\end{proof}
% END ARO12

\subsection{Operators, multiplicities, and the original good rows}

% BEGIN ARO13
\paragraph{ARO13. All event kernels are admissible orthogonal kernels.}
\label{claim:ARO13}
For all sets \(S,T\) of tuples,
\[
 e_T(S)*e_T(T)=e_T(S\cap T),\quad e_T(S)^*=e_T(S),
 \quad e_T(\varnothing)=0,\quad e_T(\mathrm{all})=e.
 \tag{ARO.15}
\]
Each kernel is locally constant, supported in \(K\), integrable,
and of positive type.  Thus \(\mathrm{EtUniv}\), \(\mathrm{EtInt}\)
and every original \(\mathrm{EtAlg}(S)\) are supplied.
\begin{proof}
By ARO10, \(\widehat e\) is a self-adjoint idempotent.  ARO6
makes every \(E_S\) central, so
\((\widehat eE_S)(\widehat eE_T)=\widehat eE_{S\cap T}\)
and \((\widehat eE_S)^*=\widehat eE_S\).  ARO5 transports these
identities to the kernels.  Their support and local constancy follow
from their finite reduction description; compact support makes
them integrable.  ARO6 gives the empty and universal events.

For positive type let \(h=e_T(S)\), let \(x_1,\ldots,x_m\) be
arbitrary finite-adelic elements, and let \(c_i\in\mathbb C\).
Using \(h*h=h\), \(h(x^{-1})=\overline{h(x)}\), and left Haar
invariance gives
\[
 \sum_{i,j}\overline{c_i}c_j h(x_i^{-1}x_j)
       =\int\left|\sum_j c_jh(z^{-1}x_j)\right|^2d\mu(z)
       \geq0.                                         \tag{ARO.16}
\]
Each product under the finite expansion is integrable, since its
support lies in a translate of a compact set.  Thus all exchanges
of the finite sum and the integral are justified.  This proves
positive type on the full group, not just on \(K\).
\end{proof}
% END ARO13

% BEGIN ARO14
\paragraph{ARO14. Bounded convolution on the given quotient Hilbert space.}
\label{claim:ARO14}
Let \(\mathcal H=X.B.I.L^2(\mu,X.B.\mathrm{fund})\) with its
original right regular action \(R\).  Every integrable finite kernel
\(h\) has a unique bounded operator \(R(h)\), characterized by
\[
 \langle u,R(h)v\rangle
       =\int h(x)\langle u,R(1,x)v\rangle\,d\mu(x),
 \qquad \|R(h)\|\leq\|h\|_1.                       \tag{ARO.17}
\]
Moreover \(R(h_1*h_2)=R(h_1)R(h_2)\) and
\(R(h^*)=R(h)^*\).
\begin{proof}
The local unitary groups are unimodular: the modular character is
the absolute determinant of the adjoint action on the Lie algebra,
and the invariant nondegenerate trace pairing makes this determinant
have absolute value one.  In the split case the same conclusion is
immediate on matrix units; in the nonsplit case the trace pairing
descends from the splitting extension.  Products and restricted
products preserve this assertion.  Thus the given Haar measure is
also right invariant and inversion invariant.

Right translation on the quotient preserves the measure specified
by the original measurable fundamental domain.  Consequently
\(R(ab)=R(a)R(b)\), \(R(1)=1\), and every \(R(a)\) is unitary.
The matrix coefficients are measurable: translate measurable
representatives and use a fundamental-domain representative,
then approximate \(L^2\) vectors by simple functions.  Their
absolute values are at most \(\|u\|\|v\|\).  Hence the right side
of (ARO.17) is a bounded sesquilinear form, conjugate linear in
\(u\), and Riesz representation supplies a bounded operator of
the indicated norm.  The same formula proves uniqueness.

For multiplication, the absolute value of the double integrand is
bounded by \(|h_1(x)h_2(y)|\|u\|\|v\|\), whose integral is
\(\|h_1\|_1\|h_2\|_1\|u\|\|v\|\).  Fubini and the change
of variable \(z=xy\) therefore give the convolution identity.
Taking the adjoint in (ARO.17), using \(R(1,x)^*=R(1,x^{-1})\),
and changing \(x\) to \(x^{-1}\) gives the involution identity.
This proves the original \(\mathrm{IsRConvF}\) specification,
including its existence and uniqueness assertions.
\end{proof}
% END ARO14

% BEGIN ARO15
\paragraph{ARO15. The literal event operators form one finite projection algebra.}
\label{claim:ARO15}
The original \(\mathrm{ReT}(X,S)\) equals \(P_S=R(e_T(S))\).
It satisfies
\[
 P_SP_T=P_{S\cap T},\quad P_S^*=P_S,\quad
 P_{\rm all}=X.B.R_e,
 \quad P_S=\sum_{t\in\mathcal T\cap S}P_{\{t\}}.
 \tag{ARO.18}
\]
Every \(P_S\) commutes with the entire infinity action and with
the original operator \(R_\phi=X.Rph\).
\begin{proof}
ARO13 supplies an integrable kernel, and ARO14 supplies its unique
operator.  Therefore the existential branch in the definition of
\(\mathrm{ReT}\) is realized; its zero fallback introduces no
alternative operator.  ARO13--ARO14 prove the product and adjoint
identities.  The universal kernel is \(e\), so uniqueness identifies
its operator with the given \(X.B.R_e\).  The finite split in ARO6,
linearity of inflation and convolution, and linearity of (ARO.17)
give the sum in (ARO.18).  Distinct singleton projectors therefore
have zero product.

The finite-adelic action \(R(1,x)\) commutes with \(R(y,1)\).
Inserting this equality in (ARO.17) proves that every \(P_S\)
commutes with the infinity action.  The defining weak integral for
\(R_\phi\) then gives \(P_SR_\phi=R_\phi P_S\): move the bounded
operator through the inner product and the scalar integral on
each side.  Thus all event and infinity projections concern the
same quotient Hilbert space and the same regular representation.
\end{proof}
% END ARO15

% BEGIN ARO16
\paragraph{ARO16. The exact infinity projector under the original gap.}
\label{claim:ARO16}
Assume the original \(\mathrm{GapHyp}(D)\).  Let \(\sigma\) be
the original Hilbert tensor product of the local representations,
let \(w=\bigotimes_vw_v\) be its specified unit vector, and put
\(m(y)=\langle w,\sigma(y)w\rangle\).  Then
\(\phi(y)=(\prod_vd_v)\overline{m(y)}\) is integrable for the
original product infinity Haar measure.  The supplied
\(T=R_\phi\) satisfies
\[
 T^*=T,\quad T^2=T,\quad
                 T R(y,1)T=m(y)T.                    \tag{ARO.19}
\]
The span of \(\sigma(y)w\), over all infinity elements \(y\),
is dense in the given tensor Hilbert space.
\begin{proof}
DC14 checks the literal ordered parameters and the original gap;
DC13 identifies every supplied local \(\mathrm{IsDS}\) with the
reconstructed model at that gap.  DC6 gives integrability of its
lowest coefficient, and DC12 gives full four-vector Schur
orthogonality with its actual Haar formal degree.  The quotient
defining the source's \(d_v=\mathrm{fdegOf}\) equals that degree:
set all four vectors equal to the given unit vector in the Schur
identity.  No additional formal-degree formula is needed here.

The product has finitely many places.  Fubini gives integrability
of \(\phi\) and the product Schur identity.  Applying that identity
to two rank-one coefficients gives, in the infinity convolution
algebra,
\(\phi^*=\phi\), \(\phi*\phi=\phi\), and the coefficient identity
whose integrated form is \(R(\phi)R(y)R(\phi)=m(y)R(\phi)\).
For the latter identity expand the two convolutions and use
\[
 \int\langle a,\sigma(x)b\rangle
             \overline{\langle a',\sigma(x)b'\rangle}\,dx
  = (\prod_vd_v)^{-1}\langle b',b\rangle\langle a,a'\rangle.
\]
Products of coefficients are integrable by Cauchy--Schwarz, and
\(\phi\in L^1\) permits the outer weak integrals.  The operator
identities follow as in ARO14.  Uniqueness of the given operator
integral identifies \(R(\phi)\) with the original \(T\).

In each local irreducible unitary representation, the closed span
of the orbit of the nonzero \(w_v\) is a nonzero invariant subspace,
hence the whole representation.  Tensor products of vectors in
these orbit spans are dense in the original tensor product, whose
defining tensor map has dense span.  They are linear combinations
of the global orbit of \(w\).  This proves the asserted cyclicity.
\end{proof}
% END ARO16

% BEGIN ARO17
\paragraph{ARO17. Evaluation identifies multiplicity maps with the actual row space.}
\label{claim:ARO17}
Under the hypotheses of ARO16, let \(Q\) be any of the event
operators of ARO15, or \(X.B.R_e\).  Let
\(\operatorname{Hom}_{G_\infty}(\sigma,\mathcal H)\) mean bounded
intertwiners, and set
\[
 \mathcal M_Q=\{Q\circ A:A\in
                    \operatorname{Hom}_{G_\infty}(\sigma,\mathcal H)\}.
\]
Then evaluation at the original \(w\) is a linear isomorphism
\[
              \mathcal M_Q\ \xrightarrow{\ A\mapsto A(w)\ }
                       \operatorname{im}(QT).        \tag{ARO.20}
\]
\begin{proof}
All members of \(\mathcal M_Q\) are intertwiners by ARO15.
Cyclicity in ARO16 proves injectivity: an intertwiner vanishing on
\(w\) vanishes on its orbit span and then, by continuity, everywhere.
Schur orthogonality in ARO16 says
\(\sigma(\phi)w=w\).  Move a bounded intertwiner through the
weak integral to get \(TA(w)=A(w)\).  Thus
\((Q A)(w)=QT A(w)\), proving inclusion in the stated image.

For the reverse inclusion first let \(z\in\operatorname{im}T\).
The rank-one compression identity (ARO.19) gives
\[
 \langle R(a,1)z,R(b,1)z\rangle
       =\|z\|^2\langle\sigma(a)w,\sigma(b)w\rangle.
 \tag{ARO.21}
\]
Indeed move \(R(a,1)\) to the other side and insert \(Tz=z\)
on both sides of the inner product, then apply (ARO.19) to
\(a^{-1}b\).  For a finite linear combination, (ARO.21) shows that
\[
 \sum_i c_i\sigma(a_i)w\longmapsto
                         \sum_i c_iR(a_i,1)z
\]
is well-defined and has norm \(\|z\|\) times the norm of its
argument.  It extends to a bounded map \(A_z\) on the full Hilbert
space, including the case \(z=0\).  Translation of the dense orbit
span proves equivariance; continuity extends it to all vectors.
One has \(A_z(w)=z\).  Given \(u=QT y\), take \(z=Ty\) and
use \(Q A_z\); its value at \(w\) is \(u\).  This proves
surjectivity and hence (ARO.20).  The proof establishes the exact
row-evaluation statement and does not identify multiplicities with
Hilbert dimensions merely by analogy.
\end{proof}
% END ARO17

% BEGIN ARO18
\paragraph{ARO18. Finite-dimensionality at the precise retained hypothesis.}
\label{claim:ARO18}
Assume \(\mathrm{RowsFinDim}\) and \(\mathrm{GapHyp}(D)\).
Then \(H_X=\operatorname{im}(X.B.R_e\,T)\) and
\(\mathcal M_e\) are finite dimensional.  Every
\(\mathcal M_{P_S}\) is finite dimensional as well.
\begin{proof}
The original \(\mathrm{RowsFinDim}\) quantifies over the given
Haar measure and fundamental domain, the regular representation,
a locally constant compactly supported finite kernel, the local
irreducible unitary representations, their formal degrees, unit
vectors, integrable lowest coefficients, and the specified infinity
operator integral.  The original setting supplies the measure,
fundamental-domain and representation fields; ARO13 supplies the
kernel properties and ARO16 supplies the degree and integrability
fields.  Thus this exact retained assertion applies to \(e\) and
\(T\), giving finite-dimensionality of \(H_X\).
ARO17 transfers it to \(\mathcal M_e\).

By ARO15, \(P_SX.B.R_e=P_S\).  Consequently every
\(P_S A\) is \(P_S(X.B.R_eA)\), and hence
\(\mathcal M_{P_S}\) is an image of the finite-dimensional
\(\mathcal M_e\).  This proves the last assertion for every set
\(S\).  The proof does not claim a new general automorphic
finite-multiplicity theorem: that is exactly the named retained
\(\mathrm{RowsFinDim}\) premise used in its original scope.
\end{proof}
% END ARO18

% BEGIN ARO19
\paragraph{ARO19. Additivity of the original row multiplicities.}
\label{claim:ARO19}
Under the hypotheses of ARO18, the original quantity
\(d(S)=\mathrm{dimX}(X,S)\) satisfies
\[
                  d(S)=\sum_{t\in\mathcal T\cap S}d(\{t\}).
 \tag{ARO.22}
\]
This proves the entire \(\mathrm{RowsAdd}(X)\).
\begin{proof}
By definition \(d(S)=\dim_{\mathbb C}\mathcal M_{P_S}\).
For an intertwiner \(A\), ARO15 decomposes \(P_SA\) as the
finite sum of \(P_{\{t\}}A\) over the indicated labels.  Conversely,
if \(A_t\in\mathcal M_{P_{\{t\}}}\), then
\(P_S A_t=A_t\); hence any finite sum belongs to
\(\mathcal M_{P_S}\).  The sum is direct: applying
\(P_{\{u\}}\) to a vanishing sum singles out its \(u\)-summand,
because the singleton projectors are pairwise orthogonal.
All spaces are finite dimensional by ARO18, so dimensions add.
The finite set \(\mathcal T\) is exactly the source's
\(\mathrm{labT}\), by the coefficient bijection and invalid-label
convention in ARO6.  This proves (ARO.22) for all sets, not only
for cylinders or finite subsets chosen in advance.
\end{proof}
% END ARO19

% BEGIN ARO20
\paragraph{ARO20. Good singleton ranges and their exact dimensions.}
\label{claim:ARO20}
Under the hypotheses of ARO18, let \(G\) be the original good operator,
namely convolution by
\(\operatorname{inf}(E_{\rm good})\), and put
\[
 V_t=H_X\cap\ker(P_{\{t\}}-1),\qquad
 H_{\rm good}=H_X\cap\ker(G-1).
\]
Then \(V_t\perp V_u\) for \(t\ne u\);
\(V_t\subseteq H_{\rm good}\) for every good tuple \(t\); and
\[
                         d(\{t\})=\dim_{\mathbb C}V_t.
 \tag{ARO.23}
\]
\begin{proof}
The good operator exists by ARO4 and ARO14.  Uniqueness identifies
it with any supplied \(\mathrm{GoodOp}\).  On the finite side,
\(E_{\rm good}(\widehat eE_{\{t\}})=\widehat eE_{\{t\}}\)
for a good \(t\).  ARO5 and ARO14 give
\(GP_{\{t\}}=P_{\{t\}}\), proving the good-range inclusion.
If \(v=P_{\{t\}}v\) and \(u=P_{\{u\}}u\) with distinct labels,
self-adjointness and zero product in ARO15 give \(\langle v,u\rangle=0\).

To prove the dimension assertion, identify the spaces exactly.
Since \(P_{\{t\}}\) commutes with \(T\), is idempotent, and is
absorbed by \(X.B.R_e\), its image after \(T\) lies in \(V_t\).
Conversely, if \(v=X.B.R_eTy\) and \(P_{\{t\}}v=v\), then
\(v=P_{\{t\}}Ty\).  Hence
\(\operatorname{im}(P_{\{t\}}T)=V_t\).
ARO17 identifies this image with \(\mathcal M_{P_{\{t\}}}\),
whose dimension is the original \(d(\{t\})\).  This proves (ARO.23)
for every valid tuple; no abstract replacement for the original
multiplicity space has been inserted.
\end{proof}
% END ARO20

% BEGIN ARO21
\paragraph{ARO21. The given good rows have exactly the singleton multiplicities.}
\label{claim:ARO21}
Retain the original isometry
\(\iota:D.s.HH\simeq H_X\), its \(\mathrm{RowDictTie}\) with
the good operator, and its \(\mathrm{LabTie}\).  Let
\(b_\sigma=\iota(D.sEE(D.\mathrm{what}_2.\mathrm{onb}_\sigma))\)
and let \(t_\sigma\) be the original tuple of labels of that row.
For every good valid tuple \(t\),
\[
                   \#\{\sigma:t_\sigma=t\}=d(\{t\}).
 \tag{ARO.24}
\]
Together with the label and goodness clauses this is
\(\mathrm{RowsGood}(X)\).
\begin{proof}
The literal row-dictionary tie is an isometric linear isomorphism
\(j\) onto the conjugate good space, followed by complex
conjugation.  Thus \(b_\sigma=\overline{j(\mathrm{onb}_\sigma)}\).
These vectors are linearly independent: conjugate a vanishing
linear relation, use independence of the original orthonormal
basis and injectivity of \(j\), then conjugate the coefficients
back.  They span \(H_{\rm good}\) for the same reason, using
surjectivity onto the conjugate good space.  This argument respects
the semilinear map; it does not treat conjugation as a complex-linear
identification.

The original \(\mathrm{LabTie}\) states that each \(t_\sigma\)
is valid and good and that \(b_\sigma\in V_{t_\sigma}\).
For a good \(t\), ARO20 puts \(V_t\) inside the span of all these
vectors.  If \(v\in V_t\), express it in this finite basis and
split the sum into labels equal to \(t\) and unequal to \(t\).
The first sum belongs to \(V_t\); subtracting it shows that the
second also belongs to \(V_t\).  By ARO20 every summand of the
second is orthogonal to \(V_t\), so its norm is zero.  It vanishes.
Thus the vectors with label \(t\) form a basis of \(V_t\).
ARO20 now gives (ARO.24).  Repeated labels are allowed: their
number is exactly the corresponding dimension.
\end{proof}
% END ARO21

% BEGIN ARO22
\paragraph{ARO22. Completing the label list while preserving every given good row.}
\label{claim:ARO22}
Let a finite label set \(\mathcal T\) have multiplicities
\(m_t\in\mathbb N\), a good-label predicate, and a given finite
row set \(I\) with good labels \(t_i\).  Suppose each good label
occurs in \(I\) exactly \(m_t\) times.  If
\(M=\sum_tm_t\), then there are labels on \(\{0,\ldots,M-1\}\)
and an injection from \(I\) into that set such that the original
labels are preserved, the good indices are exactly the image, and
every event \(S\) has exactly \(\sum_{t\in S}m_t\) indices.
\begin{proof}
Form the finite disjoint union
\[
                \mathcal R=I\ \sqcup\!
                  \coprod_{t\text{ not good}}\{t\}\times\{0,\ldots,m_t-1\}.
 \tag{ARO.25}
\]
Label an original row by its given label and a new pair by \(t\).
The multiplicity assumption on good labels shows
\(|\mathcal R|=\sum_tm_t=M\).  Choose any bijection between
\(\mathcal R\) and the standard \(M\)-element set.  Transport the
labels and the inclusion of the first summand through this one
bijection.  Its restriction to \(I\) is injective and preserves
labels; every added row has a nongood label and every original
row has a good label.  Counting an event in the two disjoint
summands gives the asserted formula.  This construction also works
when a multiplicity, \(I\), or the entire row space is zero.
It adds precisely the missing rows and makes no separate choice
of an enumeration for different events.
\end{proof}
% END ARO22

% BEGIN ARO23
\paragraph{ARO23. All fields of the original row-label record.}
\label{claim:ARO23}
Assume the hypotheses of ARO18 and ARO21 and the original
\(\mathrm{AutTie}(D,X.B,X.pl)\).  Then there is one label map
\(\mathrm{labAll}\) and one injection
\(\mathrm{inc}:\mathrm{Fin}(D_{\rm rows})\hookrightarrow
\mathrm{Fin}(D_{\rm all})\) such that
\[
 d(S)=\#\{\tau:\mathrm{labAll}(\tau)\in S\},
 \quad \mathrm{labAll}(\mathrm{inc}(\sigma))=t_\sigma,
 \quad \mathrm{labAll}(\tau)\text{ good}\ \Longleftrightarrow\
                           \tau\in\operatorname{im}(\mathrm{inc}).
 \tag{ARO.26}
\]
Every label is valid.  These are all fields of the original
\(\mathrm{RowsLab}(X)\).
\begin{proof}
By ARO15 the universal event operator is the original \(R_e\),
so \(d(\mathrm{all})=X.B.D\).  The dimension field of the given
\(\mathrm{AutTie}\) identifies this number with the original
\(D_{\rm all}\).  ARO19 therefore gives
\(D_{\rm all}=\sum_{t\in\mathcal T}d(\{t\})\).
ARO21 identifies every good-label fiber of the original
\(D_{\rm rows}\) rows with its required multiplicity.
Apply ARO22 with \(m_t=d(\{t\})\) and this exact row set.
Its event-count formula and ARO19 give the first identity in
(ARO.26); its inclusion and goodness statements give the other
two.  By construction labels belong to \(\mathcal T\), so they
satisfy every original \(\mathrm{IsLabel}\) clause.  The maps are
chosen once before \(S\), \(j\), or a cylinder is specified.
\end{proof}
% END ARO23

% BEGIN ARO24
\paragraph{ARO24. One actual AllRows object with all label and bundle ties.}
\label{claim:ARO24}
Under the hypotheses of ARO23 there is an original
\(A:\mathrm{AllRows}(D)\) with
\(\mathrm{LabelCountGoodInjective}(A)\) and
\(\mathrm{AllJoinA}(A)\).  It has \(A.B=X.B\), the event kernels
and operators of ARO12--ARO15, and
\[
                 A.eVec(j,v)=e_J(q_j,n,k_0,s,\psi_j^0)
                    \quad\hbox{for every }j,v,        \tag{ARO.27}
\]
where one nontrivial \(\psi_j^0\) is fixed for each prime.
\begin{proof}
Take, for example, \(\psi_j^0(a)=\exp(2\pi i\widetilde a/q_j)\).
This is nontrivial since \(q_j>1\).  The local idempotence in
(ARO.27) is the normalized subgroup-character identity of RRT1.
Use the label map and injection of ARO23, and use the given
\(X.B\).  The dimension identities \(A.B.I.n=n\) and
\(A.B.I.f=f\) are the original rank field and the degree equality
under \(X.\iota_E\).  ARO14--ARO15 give every operator specification
and the universal operator; ARO13 gives the universal kernel;
ARO12 gives the full cylinder identity.  These are every field
of the \(\mathrm{AllRows}\) structure.  ARO23 supplies all four
clauses of \(\mathrm{LabelCountGoodInjective}\).

For \(\mathrm{AllJoinA}\), its level clause is (ARO.27) with the
chosen nontrivial characters.  Its bundle clause is the original
\(\mathrm{RowsTie}\): use \(X.pl\) and the given \(\mathrm{AutTie}\),
the original Harish--Chandra parameter identities in that tie, and
complete splitting of \(q_0=X.B.g_0\) transported by
\(X.\iota_E\).  Complete splitting follows from the original
full list of residue-degree-one places; the sum of residue degrees
is the field degree.  Hence no extra ramification can occur.
All these fields refer to the same original bundle.  In particular
the row labels have not been supplied by an unrelated finite model.
\end{proof}
% END ARO24

\subsection{The count-idempotent interface and its common-family scope}

% BEGIN ARO25
\paragraph{ARO25. Partner-prime unipotence and the full admissibility fields.}
\label{claim:ARO25}
For every \(j\in\{0,1\}\) and every set \(S\) of its tuples,
\(e_T(\operatorname{cyl}(j,S))\) is an original admissible
idempotent at level \(g\), with distinguished unipotence prime
\(q_{1-j}\).
\begin{proof}
Decompose \(\mathcal G=\mathcal G_{1-j}\times\mathcal G_j\).
ARO6 shows that the cylinder projector is \(1\otimes Z_S\),
where \(Z_S\) is the central event projector in the \(j\)-factor.
ARO10 gives \(\widehat e=a_{1-j}\otimes a_j\), where each
\(a_i\) is the product of the corresponding conjugated character
idempotents.  Thus
\[
       \widehat{e_T(\operatorname{cyl}(j,S))}
                         =a_{1-j}\otimes(a_jZ_S).     \tag{ARO.28}
\]
If a coefficient is nonzero, its \((1-j)\)-part lies in the
product of the conjugated \(J\)'s.  Its reduction there is
unipotent.  The conjugate reductions are inverse transposes and
also unipotent.  Thus \((x-1)^n\widehat L\subseteq
q_{1-j}\widehat L\), which is the original unipotence-modulo-prime
condition.  The same conclusion follows by lifting that one
factor to \(X.K_2\) using ARO2 and ARO8 and comparing modulo
\(K(q_{1-j})\) by ARO1.

ARO13 supplies local constancy, convolution idempotence,
self-adjointness, positive type and support in \(K\).
ARO4--ARO5 supply invariance under multiplication on both sides
by \(K(g)\).  Together with the preceding unipotence these are
every field of the original \(\mathrm{IsAdmIdem}\).  They also
hold when the cylinder projector is zero.  The prime used for
unipotence is deliberately the complementary one, exactly as in
the original count application.
\end{proof}
% END ARO25

% BEGIN ARO26
\paragraph{ARO26. Every cylinder is an original counting idempotent.}
\label{claim:ARO26}
For every \(j,S\), the kernel of ARO25 satisfies the original
\(\mathrm{IsCountIdem}(\mu,q_{1-j},e_T(\operatorname{cyl}(j,S)))\).
Consequently the common object of ARO24 has the full
\(\mathrm{PCountTie}\), for every set \(S\).
\begin{proof}
Use exponents \(a=b=1\), second prime \(p=q_j\), and the
reduction datum
\[
 K\longrightarrow\mathcal G_{1-j}\times\mathcal G_j
\]
from ARO1--ARO2.  It is multiplicative and surjective; equality
of first, respectively second, coordinates is equivalent to
membership of the quotient \(x^{-1}y\) in \(K(q_{1-j})\),
respectively \(K(q_j)\).  Extending the reduction map by the
identity off \(K\) gives exactly the original total-map convention.

Set \(f_0=a_{1-j}\), \(\varepsilon_2=a_j\), and \(z=Z_S\)
in (ARO.28).  Each \(a_i\) is the character idempotent of the
product subgroup \(\prod_vP_{i,v}J_{q_i}P_{i,v}^{-1}\): its
character is the product of the local characters in ARO9, is
multiplicative there, and has modulus one.  Its coefficients are
the inverse character divided by the order of that subgroup.
Thus both \(f_0\) and \(\varepsilon_2\) satisfy all the original
\(\mathrm{IsCharIdemFin}\) clauses.  ARO6 makes \(z\) a central
idempotent.  ARO25 proves the required unipotence when the first
coefficient is nonzero.

Finally ARO1 and elementary coprime-lattice congruence give
\(K(q_0)\cap K(q_1)=K(g)\).  Equations (ARO.5) and (ARO.28)
therefore give exactly the required formula
\[
 e_T(\operatorname{cyl}(j,S))(x)=
 \begin{cases}
 \displaystyle\frac{f_0(\rho_{1-j}(x))
                 (z\varepsilon_2)(\rho_j(x))}
                 {\mu(K(q_{1-j})\cap K(q_j))},&x\in K,\\[6pt]
 0,&x\notin K.
 \end{cases}                                         \tag{ARO.29}
\]
The denominator is the original \(V\), not either one-prime
volume.  The two primes are distinct by the datum.  This supplies
every existential witness and every clause of
\(\mathrm{IsCountIdem}\).  The right disjunct of the literal
\(\mathrm{PCount}\) now gives \(\mathrm{PCountTie}\).
Arbitrary \(S\) is permitted because ARO6 reduces it to its
finite set of valid labels.
\end{proof}
% END ARO26

% BEGIN ARO27
\paragraph{ARO27. The exact common count parameters on the original bundle.}
\label{claim:ARO27}
Retain \(\mathrm{AutTie}\), an original size datum \(L\), the
original \(\mathrm{CountTheoremHypotheses}(D,L,n_0,A,C,c_b)\),
and \(4n<q_1\).  Then the full original
\(\mathrm{CountBase}(X,n_0,A,C,c_b)\) holds with
\[
 b=L.sc.b,\qquad h=Cn^{19}f^5\log(b\,\mathrm{rd}_E).
 \tag{ARO.30}
\]
The same two numbers work for both auxiliary primes and every
cylinder.
\begin{proof}
The retained count hypotheses assert the original discriminant
identity, \(n_0\leq n\), \((s+1)^2\leq n\),
\(\log\log\mathrm{rd}_E\leq n\),
\(|\lambda_{\tau,a}|\leq b^2\), \(g\leq b^A\),
\(4n<q_0\), complete splitting of \(q_0\), the gap at least
the value \(h\) in (ARO.30), and
\[
 \log f\leq\sqrt n,\qquad 2\leq b,\qquad
              \log\log b\leq c_bnf\log n.            \tag{ARO.31}
\]
Their discriminant clause is retained although this particular
transport does not need to re-use it.  The given \(\mathrm{AutTie}\)
identifies the bundle's rank, real degree, root discriminant,
ordered parameters and signs with these exact quantities.  In
particular, choose \(p\) and \(q\) to be the numbers of true and
false entries of the original \(\mathrm{epsS}\).  The
\(\mathrm{eps\_eq}\) field gives the same
\(p\mathrm{At}(v)=p\) and \(q\mathrm{At}(v)=q\) at every
infinity place, as required by \(\mathrm{SetBase}\).  The
number of negative signs is at most \(s+1\), so the signature
square bound follows from \((s+1)^2\leq n\).  The original
Harish--Chandra parameter conditions transport with the ordered
parameter; no reordering is made.

The group level is the given \(g=q_0q_1\geq1\).  Complete
splitting at both primes follows from the original datum as in
ARO24, and the two inequalities \(4n<q_0,q_1\) provide the
distinguished-prime guard for either choice of \(j\).
The gap, its lower bound by \(Cn^{19}f^5\log(b\mathrm{rd}_E)\),
and (ARO.31) then hold on the original bundle.  Finally the
\(D_*\)-field of \(\mathrm{AutTie}\) identifies
\(X.B.I.D_*(e)\) with the original positive \(D_*\).
These are all the clauses of \(\mathrm{CountBase}\), with the
single pair \((b,h)\) of (ARO.30).
\end{proof}
% END ARO27

% BEGIN ARO28
\paragraph{ARO28. Installation of the full cylinder count from its exact theorem.}
\label{claim:ARO28}
In addition to the hypotheses of ARO24 and ARO27, supply the
original assertion
\(\mathrm{CountIdempotentsClass}(\mathrm{PCount},D,n_0,A,C,c_3,c_b)\).
For any \(0<c\leq c_3\), the same \(A\) of ARO24 satisfies
\[
        \mathrm{CountAllRows}\bigl(A,\exp(-cnf\log n)\bigr).
 \tag{ARO.32}
\]
This is conditional on that exact arithmetic counting theorem;
the theorem itself is not declared proved here.
\begin{proof}
For each \(j,S\), ARO25 supplies its actual admissible kernel,
ARO26 supplies the \(\mathrm{PCount}\) predicate, and ARO27
supplies the same parameters \(b,h\).  The discrete-series,
unit-vector, Hilbert tensor, fundamental-domain, and regular-action
fields are those of \(X\), with the operator specification of
ARO14--ARO15.  Thus every argument of the supplied counting
theorem is available on this very bundle.  Its conclusion is
\[
 \left|d(\operatorname{cyl}(j,S))-
                D_*^A(e_T(\operatorname{cyl}(j,S)))\right|
 \leq e^{-c_3nf\log n}\,
                D_*^A(e_T(\operatorname{cyl}(j,S))).   \tag{ARO.33}
\]
This proves the original \(\mathrm{CountCylAt}\) on all cylinders.
Here \(n\geq2\) follows from the original gap hypothesis and
\(f\geq1\) from the degree.  Therefore
\(0<e^{-cnf\log n}<1\).  The left side of (ARO.33) is
nonnegative and its exponential factor is positive, so its right
normalizing \(D_*^A\) is nonnegative, including for zero kernels.
Since \(c\leq c_3\), replacing the exponential factor by
\(e^{-cnf\log n}\) preserves the inequality.  ARO27 supplies
\(D_*^A(e)>0\).  These are precisely all four clauses of
\(\mathrm{CountAllRows}\).  Neither a new row object nor a
different measure or formal degree is introduced in this step.
\end{proof}
% END ARO28

% BEGIN ARO29
\paragraph{ARO29. The completed local weights apply to this same row object.}
\label{claim:ARO29}
Retain the original normalized-type compatibility, admissibility
\(\mathrm{TieAdmG}\), and range \(\mathrm{RangeW}\), together
with the actual-rank threshold of FLC19.  Then the object of ARO24
satisfies the full \(\mathrm{WeightAllRows}\), with exactly the
original \(W_g\).  Its universal and two finite bad-cylinder
normalizations are
\[
 D_*^A(e)=M>0,\qquad
 D_*^A(e_T(\operatorname{cyl}(j,S_j)))=Mp_j,
 \quad p_j=\operatorname{locShare}(A,j,S_j)\geq0.      \tag{ARO.34}
\]
\begin{proof}
ARO24 provides the original \(\mathrm{AllJoinA}\), with its
precise \(e_J\) vectors and one chosen nontrivial character at
each prime.  Thus FLC19 applies with all of its original finite
type, range and admissibility inputs.  It supplies the unchanged
\(W_g\) estimate and nonnegativity for every local label and
every finite bad event.  In the selected eventual family its
additional absolute rank threshold is absorbed exactly as in FLC19,
using the original \(x^2\leq f\leq n\) growth.

For \(S_j\), take the finite image of the bad labels among the
all-row labels of ARO23.  Membership in the cylinder is exactly
badness of the row at prime \(j\); multiplicities are counted by
(ARO.26).  This is the finite-event construction in GQT2.
ARO12 gives \(e_T(\operatorname{cyl}(j,S_j))(1)=e(1)p_j\),
with real values and nonnegative \(p_j\).  By definition \(D_*^A(h)\)
is the same product of formal degree and quotient volume times
\(\operatorname{Re}h(1)\), for every \(h\).  Hence it gives
(ARO.34).  Its positivity for the full event is the original
\(\mathrm{AutTie}\) field used in ARO27.  In particular the
finite weight theorem and the arithmetic count theorem now refer
to exactly the same operators and labels.
\end{proof}
% END ARO29

% BEGIN ARO30
\paragraph{ARO30. Exact common-family consumer and the remaining existence inputs.}
\label{claim:ARO30}
Fix the original family \(\mathcal F\), its \(0<\epsilon\leq1\),
\(\epsilon_z=\epsilon/J\) with \(J\geq1\), and one prior
\(c>0\).  Suppose one threshold precedes every varying member
\(D=\mathcal F.D(x,i)\), and that each such member supplies one
\(X,G,\iota\) with the original \(\mathrm{AutTie}\),
\(\mathrm{RowDictTie}\) and \(\mathrm{LabTie}\), the gap and
count-parameter inputs of ARO27, the exact counting theorem of
ARO28 with \(c\leq c_3\), and the finite local hypotheses of
ARO29.  Retain the original \(\mathrm{RowsFinDim}\) and
\(\mathrm{SizesRec}(x,\epsilon_z,C_s)\).
Then the literal \(\mathrm{LinkIRR}(\mathcal F)\) follows.
\begin{proof}
At each member apply ARO24 once to its given \(X,G,\iota\).
ARO28 gives \(\mathrm{CountAllRows}\) for this object, and ARO29
gives \(\mathrm{WeightAllRows}\) for the same object.  ARO23
gives the exact original label injection and good-image condition.
Thus all inputs of GQT5 are supplied, with the original dimension
\(D_{\rm all}\), retained dimension \(D_{\rm rows}\), and stored
defect \(D.\mathrm{what}_2.\mathrm{delp}\).
Take the maximum of the common input threshold, the actual-rank
threshold in ARO29, and the numerical thresholds already proved
in GQT5.  None depends on the varying index \(i\).  GQT5 then
gives, for every later member, the original nonnegative defect and
its bound \(\exp(-x^{\epsilon_z/7})\), which are exactly
\(\mathrm{LinkIRR}\).

The same reasoning only needs the three arithmetic inequalities
for the full set and the two finite bad cylinders in GQT2,
provided their relative error is uniformly at most \(1/2\).
Their exact normalizing terms are now proved in ARO29.  This
weaker alternative does not assert that those three estimates
hold without an arithmetic count proof.
\end{proof}
% END ARO30

\paragraph{Scope of the completed producer.}
The original reduction, Haar normalization, level-character shape,
\(\mathrm{EtOne}\), event-operator algebra, \(\mathrm{RowsAdd}\),
\(\mathrm{RowsGood}\), \(\mathrm{RowsLab}\), common
\(\mathrm{AllRows}\) constructor and \(\mathrm{PCountTie}\) have
human proofs above.  The full pointwise count installation is proved
with the original counting theorem as an explicit antecedent.
The remaining global work includes producing the original common
\(X,G,\iota\) and all three ties in the selected family, and
supplying that arithmetic counting theorem or its three sufficient
special cases.  Original \(\mathrm{RowsFinDim}\) remains explicit.
Consequently this chapter does not by itself close the selected
\(\mathrm{GLOBAL\_RETAINED\_DEFECT}\) obligation.  Its contribution
is that these remaining inputs now feed a completely specified
mathematical row construction, with every normalization and
multiplicity retained.

\paragraph{Primary-source map.}
The attached metadata binds every claim to exact bytes of commit
\texttt{7cbf7a0c32fca1badf05931879a37c3f8a2b03b4}.
The principal source contracts are \path{ArithmeticCylinders.lean}
(\(\mathrm{AllRows}\)), \path{AllRowsOf.lean},
\path{AllRowsRhoX.lean}, \path{RhoSurjPf.lean},
\path{KppFwdPf.lean}, \path{AllRowsEtOne.lean},
\path{AllRowsAdd.lean}, \path{AllRowsDcount.lean},
\path{RowsGoodRed.lean}, \path{AllRowsLab.lean},
\path{AllRowsPCount.lean}, and \path{AllRowsAtRecord.lean}.
The direct reverse flag argument in ARO8 and the elementary trace
test in ARO9 simplify intermediate source bookkeeping; their full
original conclusions and original data are unchanged.
